% !TeX program = lualatex
\documentclass[11pt]{article}
\usepackage[margin=25mm]{geometry}
\usepackage{fontspec}
\setmainfont{Latin Modern Roman}
\usepackage{amsmath,amssymb,mathtools}
\usepackage{unicode-math}
\setmathfont{Latin Modern Math}
\usepackage{xurl,hyperref}
\hypersetup{colorlinks=true,urlcolor=blue}
\setcounter{secnumdepth}{0}
\setlength{\parindent}{0pt}
\setlength{\parskip}{5pt}
\setlength{\emergencystretch}{3em}
\newcommand{\R}{\mathbb R}
\newcommand{\N}{\mathbb N}
\newcommand{\Z}{\mathbb Z}
\newcommand{\Q}{\mathbb Q}
\newcommand{\C}{\mathbb C}
\newcommand{\D}{\mathbb D}
\newcommand{\F}{\mathbb F}
\newcommand{\sref}[2]{\hyperlink{source:#1:#2}{[S#2]}}
\newcommand{\cataloguescope}[1]{#1}
\begin{document}
\section*{Joint SLE and CLE limits at FK weights from one through four}\label{statement}
\subsection{Definitions and mathematical statement}
\paragraph{Statement recorded in the supplied source.}
For every $q\in [1,4]$, do critical square-lattice FK Dobrushin exploration curves and their nested loop collections converge jointly along every admissible marked Jordan-domain mesh approximation to chordal $SLE_\kappa (q)$ and the corresponding CLE decoration, including $q=4$?
\subsection{Short English statement}
For every $q\in [1,4]$, do critical square-lattice FK Dobrushin exploration curves and their nested loop collections converge jointly along every admissible marked Jordan-domain mesh approximation to chordal $SLE_\kappa (q)$ and the corresponding CLE decoration, including $q=4$?
\subsection{Sources}
\begin{itemize}
\item[S1] Original question source or mathematical context.
\url{https://github.com/openai/math/blob/adc7f1241b42e322a6451854ab7e4b4c146bf78a/preprints/Conformal-Limits-of-Critical-Square-Lattice-Random-Cluster-Interfaces-October-5-2026/build/sections/introduction.tex}.
\item[S2] OpenAI, manuscript citation page. Source of the formulation motivating this local problem.
\url{https://github.com/openai/math/blob/adc7f1241b42e322a6451854ab7e4b4c146bf78a/preprints/Conformal-Limits-of-Critical-Square-Lattice-Random-Cluster-Interfaces-October-5-2026/README.md}.
\item[S3] OpenAI, manuscript citation page. Source of the formulation motivating this local problem.
\url{https://github.com/openai/math/blob/adc7f1241b42e322a6451854ab7e4b4c146bf78a/preprints/Square-lattice-FK-interfaces-and-nested-loops-for-1-leq-q-lt-4-September-23-2026/README.md}.
\end{itemize}
\end{document}
